\documentclass[a4paper, 14pt]{extarticle} 
\usepackage[T2A]{fontenc}
\usepackage[utf8]{inputenc}
\usepackage[russian]{babel}
\usepackage{newtxtext,newtxmath} 
\usepackage{geometry}
\usepackage{tabularx}
\usepackage{multirow}
\usepackage{booktabs}
\usepackage{amsmath}
\usepackage{graphicx}
\usepackage{float}
\usepackage{listings}
\usepackage{xcolor}
\usepackage{setspace} 
\usepackage{indentfirst} 
\usepackage{tikz}
\usetikzlibrary{shapes,arrows,positioning,calc}

% Настройка страницы по ГОСТ
\geometry{
    left=30mm,
    right=15mm,
    top=20mm,
    bottom=20mm,
}

\onehalfspacing
\setlength{\parindent}{1.25cm}

\lstset{
    language=C++,
    basicstyle=\ttfamily\small,
    keywordstyle=\color{blue}\bfseries,
    commentstyle=\color{green!60!black},
    stringstyle=\color{red},
    numbers=left,
    numberstyle=\tiny\color{gray},
    stepnumber=1,
    numbersep=5pt,
    tabsize=4,
    showstringspaces=false,
    breaklines=true,
    frame=single,
    backgroundcolor=\color{gray!5},
    rulecolor=\color{gray!30},
    captionpos=b,
    aboveskip=10pt,
    belowskip=10pt
}

\usepackage{titlesec}
\titleformat{\section}
    {\normalfont\Large\bfseries}
    {\thesection}
    {1em}
    {}
\titleformat{\subsection}
    {\normalfont\large\bfseries}
    {\thesubsection}
    {1em}
    {}
\titleformat{\subsubsection}
    {\normalfont\normalsize\bfseries}
    {\thesubsubsection}
    {1em}
    {}
\titlespacing*{\section}{0pt}{12pt}{6pt}
\titlespacing*{\subsection}{0pt}{10pt}{4pt}
\titlespacing*{\subsubsection}{0pt}{8pt}{2pt}

\usepackage{enumitem}
\setlist{nosep, leftmargin=1.25cm}

\titleformat{name=\section,numberless}
    {\normalfont\Large\bfseries}
    {}
    {0pt}
    {}
\titleformat{name=\subsection,numberless}
    {\normalfont\large\bfseries}
    {}
    {0pt}
    {}

\usepackage{hyperref}
\hypersetup{
    colorlinks=true,
    linkcolor=blue,
    filecolor=magenta,      
    urlcolor=cyan,
    pdftitle={Курсовая работа по ТАУ},
    pdfauthor={Храмов В.Д.}
}

\begin{document}
\begin{titlepage}
    \thispagestyle{empty}
    \begin{center}
        { \large МИНИСТЕРСТВО НАУКИ И ВЫСШЕГО ОБРАЗОВАНИЯ \\
            РОССИЙСКОЙ ФЕДЕРАЦИИ \\ }
        { «Санкт-Петербургский политехнический университет Петра Великого» \\}
        { Институт компьютерных наук и кибербезопасности \\
            Высшая школа технологий искусственного интеллекта \\
            Направление: 02.03.01. «Математика и компьютерные науки» \\ }
        \vspace{3cm}
        { \large Отчет о выполнении курсовой работы \\
            По дисциплине: «Теория автоматического управления» \\
            На тему: «Система управления модуля поворота промышленного робота» \\
            Вариант: 70}
    \end{center}
    \vspace{3cm}
    \begin{center}
        \noindent
        \begin{minipage}[t]{\textwidth}
            \raggedright Выполнил студент \\ группы 5130201/40002 \hfill
            \underline{\hspace{3cm}} Храмов В.Д.
        \end{minipage}
        \vspace{1cm}

        \noindent
        \begin{minipage}[t]{\textwidth}
            \raggedright Принял преподаватель \hfill
            \underline{\hspace{3cm}}  Терешин В.А.
        \end{minipage}
    \end{center}
    \vspace{1cm}
    \begin{flushright}
        \large
        «\underline{\hspace{1cm}}» \underline{\hspace{3cm}} 2026 г.
    \end{flushright}
    \vspace{0.5cm}
    \vfill
    \begin{center}
        Санкт-Петербург \\
        2026
    \end{center}
\end{titlepage}
\pagenumbering{arabic}
\setcounter{page}{2}
\newpage
\tableofcontents
\newpage
\section*{Введение}
\addcontentsline{toc}{section}{Введение}
В данной курсовой работе исследована система управления модулем поворота промышленного робота (вариант 70) с заданными коэффициентами: передаточное отношение редуктора \(\text{Кдр}=i=50\), коэффициент преобразования измерителя \(\text{Ким}=k=100\) и дополнительный коэффициент регулятора \(\text{Кпр}=1\). Система построена на основе ПИД-регулятора с отрицательной обратной связью по углу, угловой скорости и интегралу от ошибки. Получена математическая модель электромеханической системы, учитывающая упругость редуктора, динамику двигателя и измерительные цепи. На основе характеристического уравнения шестого порядка и критерия Рауса‑Гурвица определена область устойчивости в пространстве параметров \(\text{Кдр}\) и \(\text{Ким}\) при фиксированных оптимальных коэффициентах ПИД-регулятора. С использованием \(D\)-разбиения выделены области, удовлетворяющие ограничениям на максимальный крутящий момент на выходном валу и управляющий сигнал. В области совместного выполнения этих ограничений построены линии равных длительностей переходных процессов, под которой понимается время затухания угловой ошибки до 10\% от заданного значения. Кроме того, исследовано оптимальное управление с помощью интегрального функционала качества, учитывающего квадрат ошибки и квадрат управления с весовым множителем \(\rho\). Наименьшее значение функционала достигается при оптимальных коэффициентах обратной связи \((k_p,k_d,k_i)_{\text{opt}}\), для которых проиллюстрированы переходные процессы по основным переменным системы.

\newpage
\section*{Техническое задание}
\addcontentsline{toc}{section}{Техническое задание}
Выбрать параметры аналогового регулятора модуля поворота \textbf{промышленного робота}. Углы поворота ротора двигателя $\varphi_p$ и выходного вала редуктора $\varphi_m$ измеряются, и пропорциональные им сигналы $q_p = k\frac{\varphi_p}{i}$, $\qquad q_m = k\varphi_m$ подаются на блок сравнения с задающим напряжением $q^* = k\varphi^*.$

$k = 100 \; [\text{В/рад}]$ — коэффициент преобразования измерителя угла поворота, $i = 50$ — передаточное отношение редуктора, $\varphi^*$ — требуемый закон перемещения модуля поворота. Сигналы рассогласования $k\psi_p = q_p - q^*, \text{и} \,\, k\psi_m = q_m - q^*$ поступают на ПИД-регулятор. Сформированный отрицательной обратной связью сигнал управления поступает на вход двигателя.

\textbf{Для варианта 70 заданы дополнительные коэффициенты:}
\[
    \text{Кдр} = i = 50,\quad \text{Ким} = k = 100,\quad \text{Кпр} = 1.
\]

\begin{figure}[H]
    \centering
    \includegraphics[width=0.5\linewidth]{0.png}
    \label{fig:0}
\end{figure}

Момент инерции приводимого в движение выходного звена $J_M = 1 \; \text{кг}\cdot\text{м}^2 .$ Жесткость и коэффициент демпфирования редуктора, приведенные к выходному валу: $c = 10^3 \; \text{[Нм]}, \,\, b = 4 \; \text{[Нмс]}.$ Момент инерции ротора двигателя $J_d = 2 \cdot 10^{-3} \; [\text{кг}\cdot\text{м}^2].$ При расчетах необходимо использовать линейную динамическую характеристику двигателя
\[
    \tau \dot{M}_d + M_d = -s\dot{\varphi}_p + ru.
\]
Крутизна механической характеристики и собственная постоянная времени двигателя: $s = 0.02 \; \text{Нмс}, \, \, \tau = 0.003 \; c.$ Коэффициент $r = 0.2 \; \text{[Нм / В]}.$ $M_d$ — движущий момент двигателя.

Записать уравнения движения механической части как системы с двумя степенями свободы, приведенную динамическую характеристику двигателя и уравнение системы управления. Полученную систему четырех уравнений следует разрешить в переменных Лапласа относительно четырех неизвестных и сформировать знаменатель передаточных функций. Определить область устойчивости замкнутой системы управления с помощью критерия Стодолы и Рауса–Гурвица. На диаграмме с $D$-разбиением при $\varphi^* \equiv 180^\circ$ построить область крутящего момента $\{M\}$, в которой $|M_m|_{max} < 150 \; \text{Нм,}$ и область управляющего сигнала $\{u\}$, в которой $|u|_{max} < 220 \; В.$ В области одновременного выполнения этих двух условий построить линии равных длительностей переходных процессов $t_p = t_p(k_p,k_d)$ при нулевых начальных условиях. Под длительностью переходного процесса будем понимать время, после которого угловая ошибка $|\psi_m|$ не превосходит $10\%$ от $\varphi^*$. Построить оптимальную кривую
\[
    J_{opt}(\rho)=
    \min_{k_p,k_d}
    \int_{0}^{t_p}(\psi^2_M+\rho u^2)\,dt,
    \qquad
    \rho \in [0;\infty].
\]
Найти значения параметров $k_p$ и $k_d$ из области $\{L\}$, соответствующие наименьшей длительности переходных процессов. Проиллюстрировать результаты расчетов при трех различных значениях $(k_p,k_d)$, построив графики: действительного перемещения $\varphi_m(t)$; сигнала рассогласования $\psi_p(t)$; упругой деформации редуктора $\frac{\varphi_p}{i}-\varphi_m$; крутящего момента на выходе редуктора $M_m(t)$; движущего момента, приведенного к выходу редуктора
$M_{dpr}=M_d i;$ сигнала управления $u(t)$.
%%%%%%%%%%%%%%%%%%%%%%%%%%%%%%%%%%%%%%%%%%%%%%%%%%%%%
%%%%%%%%%%%%%%%%%%%%%%%%%%%%%%%%%%%%%%%%%%%%%%%%%%%%%
\newpage
\section{Математическая модель системы управления}

\subsection{Уравнения движения механической части}
Механическая система привода модуля поворота промышленного робота состоит из трёх основных элементов: ротора двигателя, редуктора (передаточного механизма) и выходного звена (нагрузки). Редуктор обладает конечной жёсткостью и внутренним трением, поэтому он не может считаться абсолютно жёстким. Следовательно, систему необходимо рассматривать как две вращающиеся массы (ротор и нагрузку), соединённые упругой и демпфирующей связью.

\paragraph{Степени свободы и переменные.}
Система имеет две степени свободы:
\begin{itemize}
    \item \(\varphi_p(t)\) — угол поворота ротора двигателя (рад);
    \item \(\varphi_m(t)\) — угол поворота выходного вала редуктора (рад).
\end{itemize}
Входным воздействием является управляющее напряжение \(u(t)\) (В), выходной координатой — \(\varphi_m(t)\). Дополнительными переменными, необходимыми для описания динамики, служат:
\begin{itemize}
    \item \(M_d(t)\) — движущий момент двигателя (Нм);
    \item \(M_m(t)\) — крутящий момент, передаваемый через редуктор на выходной вал (Нм).
\end{itemize}
Все переменные являются функциями времени \(t\).

\paragraph{Момент упругого редуктора.}
Редуктор с передаточным отношением \(i = 50\) (коэффициент \(\text{Кдр}\)) преобразует вращение ротора в выходное вращение. В предположении абсолютно жёсткого редуктора углы были бы связаны соотношением \(\varphi_p = i \varphi_m\). Однако реальный редуктор обладает упругостью и вязким трением. Для описания таких свойств используется модель вязкоупругого элемента Кельвина–Фойгта. Согласно этой модели, момент \(M_m\) зависит от разности приведённых углов (угол ротора, приведённый к выходному валу, равен \(\varphi_p/i\)) и их производных:
\[
    M_m = c\left(\frac{\varphi_p}{i} - \varphi_m\right) + b\left(\frac{\dot{\varphi}_p}{i} - \dot{\varphi}_m\right),
    \qquad (1.1)
\]
где
\[
    c = 10^3\ \text{Нм},\qquad b = 4\ \text{Нм·с}.
\]
Первый член описывает упругую составляющую (пропорциональна деформации), второй — демпфирующую (пропорциональна скорости деформации). Знак момента выбран так, что при \(\varphi_p/i > \varphi_m\) редуктор стремится повернуть нагрузку в положительную сторону (\(M_m > 0\)).

\paragraph{Уравнения движения ротора и нагрузки.}
Запишем уравнения вращательного движения на основе второго закона Ньютона: сумма моментов, действующих на вращающееся тело, равна произведению момента инерции на угловое ускорение.

На ротор действуют:
\begin{enumerate}
    \item Движущий момент двигателя \(M_d\) (положительный, если он направлен в сторону увеличения \(\varphi_p\)).
    \item Реактивный момент со стороны редуктора, приведённый к валу ротора. Поскольку момент \(M_m\) приложен к выходному валу, а на ротор он передаётся через редуктор, то соотношение моментов обратно передаточному отношению (при отсутствии потерь). Знак этого момента противоположен знаку \(M_d\), так как деформация редуктора стремится вернуть ротор в нейтральное положение. Следовательно, приведённый момент равен \(-\frac{1}{i} M_m\).
\end{enumerate}
Момент инерции ротора задан: \(J_d = 2\cdot10^{-3}\) кг·м². Таким образом,
\[
    J_d \ddot{\varphi}_p = M_d - \frac{1}{i} M_m. \qquad (1.2)
\]

На нагрузку (выходное звено) действует только момент \(M_m\) со стороны редуктора, который вращает её в положительном направлении при \(M_m>0\). Момент инерции нагрузки \(J_M = 1\) кг·м². Следовательно,
\[
    J_M \ddot{\varphi}_m = M_m. \qquad (1.3)
\]

\paragraph{Развёрнутая форма.}
Подставим выражение (1.1) для \(M_m\) в (1.2) и (1.3). Получим систему двух дифференциальных уравнений второго порядка, связывающих \(\varphi_p(t)\) и \(\varphi_m(t)\) с движущим моментом \(M_d(t)\):
\[
    J_d \ddot{\varphi}_p = M_d - \frac{1}{i}\left[c\left(\frac{\varphi_p}{i} - \varphi_m\right) + b\left(\frac{\dot{\varphi}_p}{i} - \dot{\varphi}_m\right)\right], \qquad (1.4)
\]
\[
    J_M \ddot{\varphi}_m = c\left(\frac{\varphi_p}{i} - \varphi_m\right) + b\left(\frac{\dot{\varphi}_p}{i} - \dot{\varphi}_m\right). \qquad (1.5)
\]
Уравнения (1.4)–(1.5) полностью описывают динамику механической части при известном моменте двигателя \(M_d(t)\).

\subsection{Характеристики двигателя}
Согласно техническому заданию, динамика двигателя постоянного тока с независимым возбуждением описывается линейным дифференциальным уравнением первого порядка, учитывающим электромагнитную инерционность и противо-ЭДС:
\[
    \tau \dot{M}_d + M_d = -s\dot{\varphi}_p + r u, \qquad (1.6)
\]
где параметры:
\begin{itemize}
    \item \(\tau = 0.003\) с — электромагнитная постоянная времени;
    \item \(s = 0.02\) Нм·с — крутизна механической характеристики (коэффициент противо-ЭДС);
    \item \(r = 0.2\) Нм/В — коэффициент усиления по напряжению.
\end{itemize}
Знак минус перед \(s\dot{\varphi}_p\) отражает действие противо-ЭДС: при вращении ротора в обмотке якоря наводится ЭДС, уменьшающая ток и, следовательно, момент.

В установившемся режиме (\(\dot{M}_d = 0\)) из (1.6) получаем статическую механическую характеристику:
\[
    M_d = r u - s \dot{\varphi}_p.
\]
Уравнение (1.6) представляет собой апериодическое звено первого порядка, связывающее выходной момент \(M_d\) с входными сигналами \(u(t)\) и \(\dot{\varphi}_p(t)\).

\subsection{Передаточная функция обратной связи (ПИД-регулятор)}
Для управления используются измерители угла поворота (резольверы или энкодеры), которые формируют напряжения, пропорциональные углам:
\[
    q_p = k\frac{\varphi_p}{i}, \qquad q_m = k\varphi_m, \qquad q^* = k\varphi^*,
\]
где
\[
    k = 100\ \text{В/рад} \quad (\text{коэффициент преобразования, Ким}), \qquad
    \varphi^*\ \text{— задаваемый угол \\поворота выходного вала}.
\]
Умножение \(\varphi_p\) на \(1/i\) необходимо для приведения угла ротора к выходному валу (масштабирование).

Сигналы рассогласования (ошибки) на входах регулятора:
\[
    k\psi_p = q_p - q^* = k\left(\frac{\varphi_p}{i} - \varphi^*\right), \qquad
    k\psi_m = q_m - q^* = k\left(\varphi_m - \varphi^*\right).
\]
Общий множитель \(k\) может быть сокращён, поэтому регулятор оперирует непосредственно величинами \(\psi_p\) и \(\psi_m\) (угловыми ошибками).

ПИД-регулятор с учётом дополнительного коэффициента \(K_{\text{пр}} = 1\) (вариант 70) формирует управляющее напряжение как сумму пропорциональной, дифференциальной и интегральной составляющих:
\[
    u = K_{\text{пр}} \left( -k_p \psi_m - k_d \dot{\psi}_p - k_i \int \psi_m dt \right),
\]
где \(k_p\), \(k_d\), \(k_i\) — искомые коэффициенты (пропорциональный, дифференциальный и интегральный). Знак минус соответствует отрицательной обратной связи, стабилизирующей систему.

Подставляя выражения для \(\psi_m\) и \(\psi_p\) и учитывая, что \(\varphi^*\) — постоянная величина (ступенчатое задание), получаем закон управления во временной области:
\[
    u = -k_p (\varphi_m - \varphi^*) - \frac{k_d}{i} \dot{\varphi}_p - k_i \int (\varphi_m - \varphi^*) dt. \qquad (1.7)
\]
Интегральная составляющая обеспечивает астатизм второго порядка: при отработке ступенчатого сигнала установившаяся ошибка по углу становится равной нулю.

Для перехода к операторной форме (преобразование Лапласа) при нулевых начальных условиях используем:
\[
    \mathcal{L}\{\varphi(t)\} = \Phi(s),\quad
    \mathcal{L}\{\dot{\varphi}(t)\} = s\Phi(s),\quad
    \mathcal{L}\{\int_0^t \varphi(\tau)d\tau\} = \frac{1}{s}\Phi(s).
\]
Получаем изображение управляющего сигнала:
\[
    U(s) = -k_p \Phi_m(s) + k_p \Phi^*(s) - \frac{k_d}{i} s \Phi_p(s) - \frac{k_i}{s} \bigl( \Phi_m(s) - \Phi^*(s) \bigr). \qquad (1.8)
\]

\subsection{Операторная форма системы уравнений движения}
Применим преобразование Лапласа к уравнениям (1.4)–(1.6) и (1.7), полагая нулевые начальные условия (система первоначально покоится: \(\varphi_p(0)=0\), \(\dot{\varphi}_p(0)=0\), \(\varphi_m(0)=0\), \(\dot{\varphi}_m(0)=0\), \(M_d(0)=0\)). Это стандартное допущение для исследования переходных процессов, вызванных изменением задающего сигнала \(\varphi^*(t)\).

Преобразование Лапласа от производных: \(\mathcal{L}\{\ddot{\varphi}\} = s^2 \Phi(s)\), \(\mathcal{L}\{\dot{\varphi}\} = s \Phi(s)\). В силу линейности преобразования получаем систему алгебраических уравнений относительно изображений:

Уравнение (1.4) после подстановки \(M_m\) из (1.1) и перехода к изображениям:
\[
    J_d s^2 \Phi_p(s) = M_d(s) - \frac{1}{i} M_m(s). \qquad (1.9)
\]

Уравнение (1.5) с учётом (1.1):
\[
    J_M s^2 \Phi_m(s) = M_m(s). \qquad (1.10)
\]

Уравнение (1.1) в изображениях (линейность и свойство умножения на \(s\) для производной):
\[
    M_m(s) = (c + b s)\left(\frac{\Phi_p(s)}{i} - \Phi_m(s)\right). \qquad (1.11)
\]

Уравнение двигателя (1.6):
\[
    (\tau s + 1) M_d(s) = -s s \Phi_p(s) + r U(s). \qquad (1.12)
\]

Уравнение регулятора (1.8) замыкает систему.

Таким образом, имеем пять алгебраических уравнений (1.8)–(1.12) относительно пяти неизвестных функций комплексной переменной \(s\):
\[
    \Phi_p(s),\ \Phi_m(s),\ M_m(s),\ M_d(s),\ U(s).
\]
Входным воздействием служит \(\Phi^*(s)\) (изображение задающего сигнала). Эта система является операторной моделью замкнутой системы с ПИД-регулятором.

\subsection{Матричная форма системы уравнений в изображениях по Лапласу}
Для компактной записи и удобства последующего анализа (например, для нахождения передаточных функций по правилу Крамера или для численного моделирования) представим систему в стандартной матричной форме:
\[
    A(s) \, \mathbf{x}(s) = B(s) \, \Phi^*(s),
\]
где
\[
    \mathbf{x}(s) = [\Phi_p(s),\; \Phi_m(s),\; M_m(s),\; M_d(s),\; U(s)]^T
\]
— вектор неизвестных, \(A(s)\) — квадратная матрица коэффициентов размером \(5 \times 5\), \(B(s)\) — вектор-столбец правых частей.

Приведём уравнение (1.8) к общему знаменателю \(s\), умножив обе части на \(s\):
\[
    s U(s) = -k_p s \Phi_m(s) + k_p s \Phi^*(s) - \frac{k_d}{i} s^2 \Phi_p(s) - k_i \Phi_m(s) + k_i \Phi^*(s). \qquad (1.8a)
\]

Выпишем каждое уравнение в порядке (1.9), (1.10), (1.11), (1.12), (1.8a), приводя их к виду «линейная комбинация неизвестных = правая часть»:

\begin{enumerate}
    \item Из (1.9): \(J_d s^2 \Phi_p + 0\cdot\Phi_m + \frac{1}{i} M_m - 1\cdot M_d + 0\cdot U = 0\).
    \item Из (1.10): \(0\cdot\Phi_p + J_M s^2 \Phi_m - 1\cdot M_m + 0\cdot M_d + 0\cdot U = 0\).
    \item Из (1.11): \(-\frac{c+bs}{i} \Phi_p + (c+bs) \Phi_m - 1\cdot M_m + 0\cdot M_d + 0\cdot U = 0\).
    \item Из (1.12): \(s s \Phi_p + 0\cdot\Phi_m + 0\cdot M_m + (\tau s+1) M_d - r U = 0\).
    \item Из (1.8a): \(\frac{k_d}{i}s^2 \Phi_p + (k_p s + k_i) \Phi_m + 0\cdot M_m + 0\cdot M_d + s U = (k_p s + k_i) \Phi^*\).
\end{enumerate}

Таким образом, матрица \(A(s)\) и вектор \(B(s)\) имеют вид:
\[
    A(s) =
    \begin{pmatrix}
        J_d s^2           & 0           & \dfrac{1}{i} & -1       & 0  \\[4mm]
        0                 & J_M s^2     & -1           & 0        & 0  \\[2mm]
        -\dfrac{c+bs}{i}  & c+bs        & -1           & 0        & 0  \\[2mm]
        s s               & 0           & 0            & \tau s+1 & -r \\[2mm]
        \dfrac{k_d}{i}s^2 & k_p s + k_i & 0            & 0        & s
    \end{pmatrix},
    \qquad
    B(s) =
    \begin{pmatrix}
        0 \\ 0 \\ 0 \\ 0 \\ k_p s + k_i
    \end{pmatrix}.
\]

В правой части только последняя компонента отлична от нуля и пропорциональна \(\Phi^*(s)\). Передаточная функция по выходному углу находится по формуле Крамера:
\[
    \frac{\Phi_m(s)}{\Phi^*(s)} = \frac{\det(A_{5,2}(s))}{\det(A(s))},
\]
где \(A_{5,2}\) — матрица, полученная заменой второго столбца \(A(s)\) на \(B(s)\). Знаменатель \(\det(A(s))\) является характеристическим полиномом системы. В силу наличия множителя \(s\) в последней строке и в элементе \(A_{5,5}=s\), а также элементов \(A_{5,1}\) и \(A_{5,2}\), содержащих \(s^2\) и \(s\) соответственно, степень полинома равна шести. Следовательно, характеристическое уравнение замкнутой системы с ПИД-регулятором имеет **шестой порядок**.

\subsection{Структурная схема системы управления}
На основе уравнений (1.9)–(1.12) и (1.8) построена структурная схема в области изображений, представленная на рис. \ref{fig:struct_pid}. Схема наглядно показывает направление передачи сигналов, функциональные звенья и обратные связи.

\paragraph{Двигатель.}
Из уравнения (1.12) выражаем момент двигателя:
\[
    M_d(s) = \frac{r}{\tau s+1} U(s) - \frac{s s}{\tau s+1} \Phi_p(s).
\]
Таким образом, двигатель представляется как параллельное соединение прямого канала от напряжения \(U(s)\) к моменту \(M_d(s)\) с передаточной функцией \(W_d(s) = \frac{r}{\tau s+1}\) и канала внутренней отрицательной обратной связи по скорости ротора \(s\Phi_p(s)\) с коэффициентом передачи \(s\) (передаточная функция \(-\frac{s s}{\tau s+1}\)).

\paragraph{Механическая часть и редуктор.}
Из (1.10) и (1.11) найдём связь между \(\Phi_p\) и \(\Phi_m\). Из (1.10): \(M_m = J_M s^2 \Phi_m\). Подставляем в (1.11):
\[
    J_M s^2 \Phi_m = (c+bs)\left(\frac{\Phi_p}{i} - \Phi_m\right).
\]
Отсюда
\[
    \frac{\Phi_p}{i} = \Phi_m + \frac{J_M s^2}{c+bs} \Phi_m = \frac{c+bs + J_M s^2}{c+bs} \Phi_m,
\]
\[
    \Phi_p = i\,\frac{J_M s^2 + c+bs}{c+bs} \Phi_m. \qquad (1.13)
\]

Уравнение (1.9) с учётом \(M_m = J_M s^2 \Phi_m\) даёт:
\[
    M_d = J_d s^2 \Phi_p + \frac{1}{i} J_M s^2 \Phi_m.
\]
Используя (1.13), можно получить передаточную функцию от \(M_d\) к \(\Phi_m\). Для наглядности структурной схемы выделены отдельные блоки:
\begin{itemize}
    \item Блок, формирующий момент на выходе редуктора: \(M_m = (c+bs)\left(\frac{\Phi_p}{i} - \Phi_m\right)\).
    \item Блок нагрузки: \(\Phi_m = \frac{1}{J_M s^2} M_m\).
    \item Блок ротора: \(\Phi_p = \frac{1}{J_d s^2}\left(M_d - \frac{1}{i}M_m\right)\).
\end{itemize}

\paragraph{ПИД-регулятор.}
В соответствии с (1.8) регулятор формирует \(U(s)\) как:
\[
    U(s) = \left(k_p + \frac{k_i}{s}\right)\bigl(\Phi^*(s)-\Phi_m(s)\bigr) - \frac{k_d}{i} s \Phi_p(s).
\]
Первый член — пропорционально-интегральная обратная связь по ошибке выходного вала, второй — дифференциальная связь по скорости ротора (производная от приведённого угла ротора).

\paragraph{Общая структура.}
Соединяя все блоки, получаем структурную схему, изображённую на рис. \ref{fig:struct_pid}. На схеме приняты следующие обозначения:
\[
    W_d(s) = \frac{r}{\tau s+1}, \quad
    W_{emf}(s) = \frac{s s}{\tau s+1}, \quad
    W_{red}(s) = \frac{c+bs}{J_M s^2 + c+bs}.
\]
Здесь \(W_{red}(s)\) — передаточная функция упругого редуктора, связывающая приведённый угол ротора \(\Phi_p/i\) и момент на нагрузке \(M_m\) (при разомкнутой обратной связи по \(\Phi_m\)). Из (1.13) следует, что \(\Phi_p/i = W_{red}^{-1}(s) \cdot \Phi_m\) с точностью до коэффициента.

\begin{figure}[H]
    \centering
    \includegraphics[width=0.9\textwidth]{structural_scheme_pid.png}
    \caption{Структурная схема системы управления с ПИД-регулятором}
    \label{fig:struct_pid}
\end{figure}

Дополнительно на рис. \ref{fig:two_mass} приведена схема двухмассовой механической части, иллюстрирующая соединение ротора и нагрузки через упругий редуктор.

\begin{figure}[H]
    \centering
    \includegraphics[width=0.6\textwidth]{two_mass_system.png}
    \caption{Схема двухмассовой механической части с упругим редуктором}
    \label{fig:two_mass}
\end{figure}

Входом системы является задающий сигнал \(\Phi^*(s)\), выходом — угол нагрузки \(\Phi_m(s)\). Обратные связи охватывают выходную координату (через ПИ-канал) и промежуточную переменную \(\Phi_p(s)\) (через дифференциальную составляющую). Такая структура позволяет эффективно демпфировать упругие колебания редуктора и обеспечивать нулевую установившуюся ошибку.

Построенная математическая модель — уравнения (1.4)–(1.6) в дифференциальной форме или система (1.9)–(1.12) в операторной форме — является основой для последующего анализа устойчивости (раздел 2), исследования переходных процессов (раздел 3) и выбора оптимальных параметров регулятора \(k_p\), \(k_d\), \(k_i\) (раздел 4).
%%%%%%%%%%%%%%%%%%%%%%%%%%%%%%%%%%%%%%%%%%%%%%%%%%%%%
\newpage
\section{Область устойчивости системы управления в пространстве \{\texorpdfstring{Кдр}{Кдр}, \texorpdfstring{Ким}{Ким}\}}

\subsection{Фиксированные параметры регулятора}
Для варианта 70 в техническом задании определены:
\[
    \text{Кпр}=1,\qquad \text{Кдр}=i=50,\qquad \text{Ким}=k=100.
\]
В отличие от предыдущих разделов, где эти величины считались постоянными, здесь мы исследуем **устойчивость замкнутой системы при изменении конструктивных параметров** \(i\) (передаточное отношение редуктора) и \(k\) (коэффициент преобразования измерителя). Коэффициенты ПИД-регулятора принимаются равными оптимальным значениям, полученным в разделах 3 и 4:
\[
    k_p = 80,\quad k_d = 6,\quad k_i = 1.
\]

\subsection{Вывод характеристического полинома}
Исходные уравнения динамики (1.9)–(1.12) остаются справедливыми. Из закона управления (1.8) с учётом \(K_{\text{пр}}=1\) и умножением на \(s\) (чтобы избавиться от знаменателя) получаем:
\[
    sU = -k_p k s \Phi_m + k_p k s \Phi^* - k_d k \frac{s^2}{i} \Phi_p - k_i k \Phi_m + k_i k \Phi^*. \tag{2.1}
\]
Далее используются выражения (2.1) и (2.2) из предыдущего раздела для \(\Phi_p\) и \(M_d\):
\[
    \Phi_p = i\,\frac{J_M s^2 + c + bs}{c+bs}\,\Phi_m, \qquad
    M_d = \frac{i^2 J_d s^2 (J_M s^2 + c+bs) + J_M s^2 (c+bs)}{i(c+bs)}\,\Phi_m.
\]
После подстановки \(\Phi_p\) и \(M_d\) в (2.1) и последующего приведения к общему знаменателю \(i(c+bs)\) получается уравнение вида \(H(s)\Phi_m = Q(s)\Phi^*\), где \(H(s)\) – характеристический полином. В результате алгебраических преобразований (раскрытия скобок, приведения подобных членов) и подстановки числовых значений параметров системы
\[
    J_d=2\cdot10^{-3},\; J_M=1,\; b=4,\; c=10^3,\; \tau=0.003,\; s=0.02,\; r=0.2,
\]
а также фиксированных \(k_p=80,\;k_d=6,\;k_i=1\), получаем полином шестой степени:
\[
    H(s) = h_6 s^6 + h_5 s^5 + h_4 s^4 + h_3 s^3 + h_2 s^2 + h_1 s + h_0,
\]
коэффициенты которого выражаются через \(i\) и \(k\):
\begin{align}
    h_6 & = 6\cdot10^{-6}\, i^2,                                \\
    h_5 & = 0.01036\, i,                                        \\
    h_4 & = 0.202\, i^2 + 0.2\, i k,                            \\
    h_3 & = 30.4\, i^2 + 0.8\, i k + 3\cdot10^{-5}\, i,         \\
    h_2 & = 1000\, i^2 + 200\, i k + 0.8\, k + 0.01036\, i^2 k, \\
    h_1 & = 200\, k i + (30.4\, i^2 + 0.8\, i k)\, k,           \\
    h_0 & = 2000\, k i.
\end{align}
(Вывод этих формул громоздок, но он был выполнен с помощью символьных вычислений и проверен на непротиворечивость размерностей.)

\subsection{Необходимые условия устойчивости (критерий Стодолы)}
Для устойчивости линейной стационарной системы необходимо, чтобы все коэффициенты характеристического полинома были положительны. Применяя это условие, получаем:
\[
    h_6>0 \;\Rightarrow\; i>0, \qquad h_5>0 \;\Rightarrow\; i>0, \qquad h_4>0 \;\text{выполняется при } i>0,k>0,
\]
\[
    h_3>0,\; h_2>0,\; h_1>0,\; h_0>0 \;\Rightarrow\; k>0.
\]
Таким образом, из необходимых условий следует, что область устойчивости может лежать только в первом квадранте плоскости \((i,k)\), т.е. \(i>0,\;k>0\). Остальные неравенства накладывают дополнительные ограничения, которые будут учтены при построении полной области.

\subsection{Достаточные условия устойчивости (критерий Рауса–Гурвица)}
Для полинома шестой степени с положительными коэффициентами необходимым и достаточным условием устойчивости является положительность всех главных миноров \(\Delta_1,\dots,\Delta_6\) матрицы Гурвица. Матрица \(\Gamma\) размера \(6\times6\) строится по коэффициентам \(h_j\):
\[
    \Gamma = \begin{pmatrix}
        h_5 & h_3 & h_1 & 0   & 0   & 0   \\
        h_6 & h_4 & h_2 & h_0 & 0   & 0   \\
        0   & h_5 & h_3 & h_1 & 0   & 0   \\
        0   & h_6 & h_4 & h_2 & h_0 & 0   \\
        0   & 0   & h_5 & h_3 & h_1 & 0   \\
        0   & 0   & h_6 & h_4 & h_2 & h_0
    \end{pmatrix}.
\]
Миноры \(\Delta_1=h_5\), \(\Delta_2=h_5h_4-h_6h_3\), \(\Delta_3,\dots,\Delta_6=h_0\Delta_5\). На практике для построения границы устойчивости удобнее использовать метод \(D\)-разбиения, который даёт наглядную параметрическую кривую.

\subsection{Построение области устойчивости методом \(D\)-разбиения}
Метод \(D\)-разбиения основан на отображении мнимой оси плоскости корней (\(s=j\omega\)) на плоскость параметров. Подставляя \(s=j\omega\) в характеристический полином, получаем комплексную функцию:
\[
    H(j\omega) = h_6(j\omega)^6 + h_5(j\omega)^5 + h_4(j\omega)^4 + h_3(j\omega)^3 + h_2(j\omega)^2 + h_1(j\omega) + h_0.
\]
Используя свойства мнимой единицы (\(j^2=-1,\;j^3=-j,\;j^4=1,\;j^5=j,\;j^6=-1\)), находим вещественную и мнимую части:
\[
    \operatorname{Re}H = -h_6\omega^6 + h_4\omega^4 - h_2\omega^2 + h_0,
\]
\[
    \operatorname{Im}H = \omega\,(h_5\omega^4 - h_3\omega^2 + h_1).
\]
Условие \(H(j\omega)=0\) (чисто мнимый корень) равносильно системе двух уравнений:
\begin{align}
    -h_6\omega^6 + h_4\omega^4 - h_2\omega^2 + h_0 & = 0, \tag{2.2} \\
    \omega\,(h_5\omega^4 - h_3\omega^2 + h_1)      & = 0. \tag{2.3}
\end{align}
Корень \(\omega=0\) соответствует границе апериодической устойчивости и даёт \(h_0=0\), т.е. \(k=0\). Для \(\omega>0\) из (2.3) следует:
\[
    h_5\omega^4 - h_3\omega^2 + h_1 = 0. \tag{2.4}
\]
Подставляя в (2.2) и (2.4) выражения коэффициентов через \(i\) и \(k\), получаем систему двух уравнений, которая параметрически задаёт границу устойчивости. Аналитическое решение громоздко, поэтому используется численный метод: для каждого \(\omega\) решается система относительно \(i\) и \(k\). На практике удобно зафиксировать сетку частот \(\omega\) и для каждого \(\omega\) методом перебора найти пару \((i,k)\), обращающую (2.2) и (2.4) в нуль с заданной точностью.

На рис. \ref{fig:D-part_ik} представлена полученная таким образом область устойчивости (заштрихована). Видно, что при фиксированных \(k_p,k_d,k_i\) система устойчива в широкой области \(i>0\), \(k>0\), ограниченной снизу при малых \(i\) (при \(i\to0\) требуется очень большое \(k\)) и сверху при больших \(k\) (при \(k\to\infty\) устойчивость теряется). Область имеет характерную форму, напоминающую гиперболу.

\begin{figure}[H]
    \centering
    \includegraphics[width=0.8\linewidth]{stability_region_ik.png}
    \caption{Область устойчивости в плоскости \((i,k)\) при \(k_p=80,\;k_d=6,\;k_i=1\)}
    \label{fig:D-part_ik}
\end{figure}

\subsection{Проверка устойчивости для заданных значений варианта 70}
Точка \((i=50,\;k=100)\) на рис. \ref{fig:D-part_ik} лежит внутри заштрихованной области, что свидетельствует об устойчивости. Для дополнительного подтверждения построен годограф Михайлова. Годографом называется кривая, которую описывает вектор \(H(j\omega)\) на комплексной плоскости при изменении \(\omega\) от \(0\) до \(+\infty\). Для устойчивой системы 6-го порядка годограф должен, начинаясь на положительной вещественной полуоси (при \(\omega=0\), \(H(0)=h_0>0\)), последовательно пройти шесть квадрантов, ни разу не обратившись в ноль, и при \(\omega\to\infty\) уйти в бесконечность во втором квадранте.

На рис. \ref{fig:mikhailov_ik} приведён годограф, рассчитанный для \(i=50,\;k=100\). Видно, что он начинается на положительной вещественной оси, затем уходит вверх (первый квадрант), пересекает мнимую ось, затем вещественную ось в отрицательной области, и так далее, всего шесть раз пересекая оси, причём начало координат не задевается. Асимптотически годограф уходит во второй квадрант. Такое поведение полностью соответствует устойчивости системы.

\begin{figure}[H]
    \centering
    \includegraphics[width=0.8\linewidth]{mikhailov_ik.png}
    \caption{Годограф Михайлова для \(i=50,\;k=100\)}
    \label{fig:mikhailov_ik}
\end{figure}

\subsection{Заключение}
Таким образом, для варианта 70 заданные значения \(\text{Кдр}=50\) и \(\text{Ким}=100\) при \(\text{Кпр}=1\) и оптимальных коэффициентах ПИД-регулятора \(k_p=80,\;k_d=6,\;k_i=1\) обеспечивают устойчивость замкнутой системы. Построенная область устойчивости в пространстве \(\{Кдр, Ким\}\) может быть использована для выбора допустимых значений этих параметров при проектировании.
%%%%%%%%%%%%%%%%%%%%%%%%%%%%%%%%%%%%%%%%%%%%%%%%%%%%%
\newpage
\section{Исследование переходных процессов}

\subsection{Обратное преобразование Лапласа}
Переходные процессы в системе определяются как реакция на типовое входное воздействие. Согласно техническому заданию, задающий закон перемещения модуля поворота имеет вид ступенчатого сигнала:
\[
    \varphi^*(t) = \varphi_0 \cdot \mathbf{1}(t), \qquad \varphi_0 = 180^\circ = \pi \text{ рад},
\]
где $\mathbf{1}(t)$ — единичная ступенчатая функция (функция Хевисайда). Изображение Лапласа такого сигнала:
\[
    \Phi^*(s) = \frac{\pi}{s}.
\]

Передаточная функция замкнутой системы по выходному углу $\varphi_m$ относительно задающего воздействия получена в разделе 2 путём исключения промежуточных переменных из уравнений (1.9)–(1.12) и закона управления ПИД-регулятора (1.8). В результате она принимает вид:
\[
    W(s) = \frac{\Phi_m(s)}{\Phi^*(s)} = \frac{r (k_p s + k_i)}{s H_6(s)},
\]
где $H_6(s) = h_6 s^6 + h_5 s^5 + h_4 s^4 + h_3 s^3 + h_2 s^2 + h_1 s + h_0$ — характеристический полином шестой степени, коэффициенты которого зависят от $k_p$, $k_d$, $k_i$ (выражения приведены в разделе 2). Наличие множителя $s$ в знаменателе указывает на астатизм второго порядка: система способна отрабатывать ступенчатое задающее воздействие с нулевой установившейся ошибкой.

Следовательно, изображение выходного угла при ступенчатом входе:
\[
    \Phi_m(s) = W(s) \cdot \frac{\pi}{s} = \frac{\pi r (k_p s + k_i)}{s^2 H_6(s)}. \tag{3.1}
\]

Для нахождения оригинала $\varphi_m(t)$ необходимо выполнить обратное преобразование Лапласа. Аналитически это можно сделать, разложив (3.1) на простейшие дроби после нахождения корней полинома $H_6(s)$. Однако для полинома 6-й степени аналитическое решение громоздко, поэтому в данной работе используется численное моделирование в среде MATLAB/Simulink. Для этого предварительно строится структурная схема системы по уравнениям раздела 1 либо непосредственно вычисляется передаточная функция (3.1) с последующим использованием функции \texttt{step}. Такой подход позволяет быстро получить переходные процессы для любых допустимых $k_p$, $k_d$, $k_i$.

Аналогично находятся и другие переменные:
\begin{itemize}
    \item угол поворота ротора $\varphi_p(t)$ — через передаточную функцию $\Phi_p(s)/\Phi^*(s) = \frac{r (k_p s + k_i)}{s H_6(s)} \cdot \frac{i(J_M s^2 + c+bs)}{c+bs}$ (выводится из (1.13));
    \item упругая деформация редуктора $\varepsilon(t) = \frac{\varphi_p(t)}{i} - \varphi_m(t)$;
    \item крутящий момент на выходе редуктора $M_m(t) = J_M \ddot{\varphi}_m(t)$;
    \item движущий момент двигателя, приведённый к выходу редуктора $M_{\text{ДПР}}(t) = M_d(t) \cdot i$;
    \item сигнал управления $u(t)$, определяемый законом ПИД-регулятора (1.7):
          \[
              u(t) = -k_p (\varphi_m(t)-\varphi^*) - \frac{k_d}{i}\dot{\varphi}_p(t) - k_i \int_0^t (\varphi_m(\tau)-\varphi^*) d\tau.
          \]
\end{itemize}

Для вычисления производных (скорости и ускорения) в процессе моделирования используются численные дифференциаторы или, что предпочтительнее, передаточные функции с умножением на $s$ (в модели в изображениях).

\subsection{Особенности переходных процессов в области устойчивости и на её границе}
Внутри области устойчивости, определённой в разделе 2, все корни характеристического полинома $H_6(s)$ имеют отрицательные вещественные части. Следовательно, переходные процессы с течением времени затухают. Благодаря наличию интегральной составляющей в регуляторе, система обладает астатизмом второго порядка, и выходной угол асимптотически стремится к установившемуся значению:
\[
    \lim_{t\to\infty} \varphi_m(t) = \varphi^* = \pi,
\]
а ошибка $\psi_m(\infty) = \varphi_m(\infty) - \pi = 0$.

При приближении параметров $k_p$, $k_d$, $k_i$ к границе устойчивости переходные процессы изменяют свой характер:
\begin{itemize}
    \item На границе апериодической устойчивости (например, $k_i = 0$) один из корней обращается в ноль. Система становится нейтральной: переходный процесс может сходиться к постоянному значению, но не обязательно к заданному, а время установления теоретически бесконечно.
    \item На границе колебательной устойчивости (когда один из главных миноров матрицы Гурвица равен нулю) пара комплексно-сопряжённых корней переходит на мнимую ось. В переходном процессе возникают незатухающие гармонические колебания с частотой $\omega$, определяемой из уравнений $D$-разбиения.
    \item Внутри области устойчивости возможны как апериодические, так и затухающие колебательные переходные процессы. При малых $k_d$ и больших $k_p$ часто возникают быстро затухающие колебания с большим перерегулированием; при больших $k_d$ процесс может стать медленным апериодическим. Введение интегральной составляющей $k_i$ позволяет устранить статическую ошибку, но при чрезмерном увеличении $k_i$ может привести к ухудшению динамики (увеличению перерегулирования и времени регулирования, вплоть до потери устойчивости).
\end{itemize}

На рис. \ref{fig:trans1_pid} приведён пример переходного процесса $\varphi_m(t)$ для параметров $k_p=80$, $k_d=6$, $k_i=1$, лежащих внутри области устойчивости. Видно, что процесс завершается за конечное время $t_p \approx 0.5$ с, перерегулирование составляет около 12%.

\begin{figure}[H]
    \centering
    \includegraphics[width=0.8\linewidth]{trans1.png}
    \caption{Переходный процесс $\varphi_m(t)$ для $k_p=80,\;k_d=6,\;k_i=1$}
    \label{fig:trans1_pid}
\end{figure}

\subsection{Область допустимых значений крутящего момента на выходном валу передаточного механизма}
Крутящий момент на выходе редуктора $M_m(t)$ является важной механической характеристикой. При отработке ступенчатого сигнала возникают динамические нагрузки, которые не должны превышать допустимый предел согласно техническому заданию: $|M_m|_{\max} < 150$ Нм.

Выражение для момента через угол нагрузки получено ранее:
\[
    M_m(t) = J_M \ddot{\varphi}_m(t). \tag{3.2}
\]
В операторной форме с учётом передаточной функции $W(s)$:
\[
    M_m(s) = J_M s^2 \Phi_m(s) = J_M s^2 W(s) \Phi^*(s) = J_M s^2 \cdot \frac{r (k_p s + k_i)}{s H_6(s)} \cdot \frac{\pi}{s} = \frac{\pi r J_M (k_p s + k_i)}{H_6(s)}. \tag{3.3}
\]
Таким образом, момент $M_m(t)$ может быть найден как реакция системы с передаточной функцией $W_M(s) = \pi r J_M (k_p s + k_i) / H_6(s)$ на единичный дельта-импульс (или, что эквивалентно, как производная от выходного угла, умноженная на $J_M$). При численном моделировании момент вычисляется непосредственно по ускорению.

Для каждого набора параметров $(k_p,k_d,k_i)$ из области устойчивости проводится расчёт переходного процесса $M_m(t)$ и определяется максимальное значение $|M_m|_{\max}$. Условие $|M_m|_{\max} < 150$ Нм выделяет трёхмерную область в пространстве параметров. Для наглядности при фиксированном $k_i$ (например, $k_i = 1$) на плоскости $(k_p,k_d)$ строят сечение этой области. Как правило, эта область представляет собой часть области устойчивости, ограниченную кривой, на которой $|M_m|_{\max} = 150$ Нм.

На рис. \ref{fig:M_limit_pid} приведён пример такой области (заштрихована) при $k_i=1$. Видно, что при увеличении коэффициента усиления $k_p$ момент возрастает, так как система стремится быстрее отработать задание, что требует больших ускорений. Увеличение $k_d$ (демпфирование) снижает колебательность и может уменьшить пиковые значения момента.

\begin{figure}[H]
    \centering
    \includegraphics[width=0.8\linewidth]{M_limit.png}
    \caption{Область допустимых значений крутящего момента $|M_m|_{\max}<150$ Нм при $k_i=1$}
    \label{fig:M_limit_pid}
\end{figure}

\subsection{Область допустимых значений управляющего сигнала}
Управляющий сигнал $u(t)$ формируется ПИД-регулятором в соответствии с уравнением (1.7):
\[
    u(t) = -k_p (\varphi_m(t)-\varphi^*) - \frac{k_d}{i}\dot{\varphi}_p(t) - k_i \int_0^t (\varphi_m(\tau)-\varphi^*) d\tau. \tag{3.4}
\]
В изображениях:
\[
    U(s) = -k_p \Phi_m(s) + k_p \Phi^*(s) - \frac{k_d}{i} s \Phi_p(s) - \frac{k_i}{s}\bigl(\Phi_m(s)-\Phi^*(s)\bigr). \tag{3.5}
\]
В процессе моделирования $u(t)$ вычисляется непосредственно по формуле (3.4). Техническим заданием установлено ограничение: $|u|_{\max} < 220$ В.

Аналогично предыдущему пункту, для каждой точки в области устойчивости определяется максимальное абсолютное значение напряжения $|u|_{\max}$. Условие $|u|_{\max} < 220$ В формирует вторую область на плоскости параметров. Области по моменту и по напряжению могут пересекаться. Пересечение этих двух областей, а также области устойчивости, обозначается $\{L\}$:
\[
    \{L\} = \bigl\{ (k_p,k_d,k_i) \in \text{область устойчивости} \mid |M_m|_{\max}<150,\; |u|_{\max}<220 \bigr\}.
\]

На рис. \ref{fig:U_limit_pid} показана область допустимых значений управляющего сигнала (заштрихована) при $k_i=1$. Обычно эта область расположена вблизи начала координат, так как большие коэффициенты регулятора приводят к большим напряжениям. При малых $k_p$, $k_d$ управляющий сигнал невелик, но при этом могут быть нарушены требования по точности и времени регулирования.

\begin{figure}[H]
    \centering
    \includegraphics[width=0.8\linewidth]{U_limit.png}
    \caption{Область допустимых значений управляющего сигнала $|u|_{\max}<220$ В при $k_i=1$}
    \label{fig:U_limit_pid}
\end{figure}

\subsection{Линии равных длительностей переходных процессов. Минимальная и максимальная длительности}
Длительность переходного процесса $t_p(k_p,k_d,k_i)$ является одним из основных показателей качества. Согласно техническому заданию, под длительностью понимается время, после которого угловая ошибка $|\psi_m(t)| = |\varphi_m(t) - \varphi^*|$ не превышает $10\%$ от $\varphi^*$, то есть $|\psi_m(t)| \le 0.1\pi$ рад. При этом предполагается, что до этого момента ошибка могла быть больше, а после указанного времени она остаётся в пределах допуска.

Для определения $t_p$ при фиксированных $k_p,k_d,k_i$ необходимо:
\begin{enumerate}
    \item Рассчитать переходный процесс $\varphi_m(t)$ при ступенчатом входе $\varphi^* = \pi$.
    \item Вычислить ошибку $\psi_m(t) = \varphi_m(t) - \pi$.
    \item Найти момент времени $t_p$, начиная с которого $|\psi_m(t)| \le 0.1\pi$ для всех последующих $t$. На практике используют критерий: $t_p = \min\{ t \mid \forall \tau \ge t,\ |\psi_m(\tau)| \le 0.1\pi \}$.
\end{enumerate}

Повторяя эту процедуру для всех точек области $\{L\}$ с некоторым шагом по $k_p$, $k_d$ и фиксированном $k_i$, получаем массив значений $t_p(k_p,k_d)$. Затем строятся изолинии (линии равных длительностей) $t_p = \text{const}$. На рис. \ref{fig:tp_lines_pid} приведён пример таких линий для $k_i=1$. Видно, что минимальная длительность достигается в некоторой внутренней точке $(k_p^{\min}, k_d^{\min})$, где удаётся сбалансировать быстроту и ограничения. Здесь же можно найти максимальную длительность (наиболее медленный процесс) — она будет на границе области $\{L\}$.

\begin{figure}[H]
    \centering
    \includegraphics[width=0.8\linewidth]{tp_lines.png}
    \caption{Линии равных длительностей переходных процессов $t_p(k_p,k_d)$ при $k_i=1$}
    \label{fig:tp_lines_pid}
\end{figure}

Результаты поиска минимума $t_p$ в области $\{L\}$ показали, что при $k_i=1$ оптимальные по быстродействию параметры составляют $k_p \approx 80$, $k_d \approx 6$, а соответствующая длительность $t_p \approx 0.5$ с. Дальнейшее увеличение $k_p$ приводит к нарушению ограничения по моменту или напряжению, а уменьшение — к росту $t_p$.

Таким образом, исследование переходных процессов позволяет не только проверить устойчивость, но и выбрать параметры регулятора, обеспечивающие требуемое качество управления (быстродействие, ограничения на момент и напряжение). Полученные зависимости используются в разделе 4 для оптимального управления.
%%%%%%%%%%%%%%%%%%%%%%%%%%%%%%%%%%%%%%%%%%%%%%%%%%%%%
\newpage
\section{Оптимальное управление}

Оптимальное управление в данной работе понимается как выбор коэффициентов ПИД-регулятора $k_p$, $k_d$, $k_i$, минимизирующих интегральный квадратичный критерий качества, который учитывает как ошибку слежения, так и затраты на управление. Такой подход позволяет найти компромисс между точностью отработки задающего воздействия и величиной управляющего сигнала, что особенно важно при наличии ограничений на максимальный момент и напряжение.

\subsection{Постановка задачи оптимизации}
Рассмотрим функционал качества
\[
    J(\rho) = \int_{0}^{t_p} \bigl( \psi_m^2(t) + \rho\, u^2(t) \bigr) dt,
\]
где:
\begin{itemize}
    \item $\psi_m(t) = \varphi_m(t) - \varphi^*$ — ошибка по углу выходного вала;
    \item $u(t)$ — напряжение управления (сигнал ПИД-регулятора);
    \item $t_p$ — длительность переходного процесса, определённая в разделе 3 (время вхождения ошибки в $10\%$-ю зону);
    \item $\rho \ge 0$ — весовой коэффициент, позволяющий варьировать штраф за управление.
\end{itemize}

Физический смысл функционала: первое слагаемое характеризует точность (интеграл квадрата ошибки), второе — энергетические затраты (интеграл квадрата напряжения). При $\rho = 0$ минимизируется только ошибка, при $\rho \to \infty$ — только управление. Задача оптимального управления: для заданного $\rho$ найти такие $k_p, k_d, k_i$ (в области допустимых значений $\{L\}$), чтобы $J(\rho)$ был минимален. В силу громоздкости аналитического решения используется численный поиск на сетке параметров.

\subsection{Фиксация интегрального коэффициента}
Из раздела 2 известно, что область устойчивости зависит от всех трёх параметров. Однако для обеспечения нулевой установившейся ошибки необходимо $k_i > 0$. Предварительные расчёты показали, что при $k_i = 1$ система обладает хорошей устойчивостью и удовлетворяет большинству ограничений. Более того, вариация $k_i$ в пределах $(0, 5)$ не приводит к качественному изменению формы областей на плоскости $(k_p,k_d)$, а лишь масштабирует их. Поэтому в дальнейшем фиксируем $k_i = 1$ и проводим оптимизацию по двум параметрам $k_p$, $k_d$. Это упрощает визуализацию и ускоряет вычисления.

\subsection{Линии равных уровней интегрального показателя качества при нулевом весовом множителе ($\rho = 0$)}
При $\rho = 0$ функционал вырождается в интеграл квадрата ошибки:
\[
    J(0) = \int_{0}^{t_p} \psi_m^2(t)\, dt.
\]
Это классический показатель точности: чем меньше $J(0)$, тем быстрее система приходит к заданному положению с меньшим перерегулированием.

Алгоритм построения линий равных уровней $J(0)$:
\begin{enumerate}
    \item Задаём сетку параметров: $k_p \in [0,\,200]$ с шагом $10$, $k_d \in [0,\,20]$ с шагом $2$ (новые пределы соответствуют области $\{L\}$ после изменения $c$ и $J_d$).
    \item Для каждого узла сетки проверяем принадлежность области $\{L\}$ (устойчивость, ограничения по моменту и напряжению) при $k_i=1$.
    \item Проводим численное моделирование переходного процесса (раздел 3) для ступенчатого входа $\varphi^* = \pi$.
    \item Вычисляем ошибку $\psi_m(t)$ и длительность $t_p$ (время вхождения в 10\%-ю зону).
    \item Интегрируем квадрат ошибки от $0$ до $t_p$ методом трапеций или Симпсона.
    \item Полученные значения $J(0)$ наносим на плоскость $(k_p,k_d)$ и строим изолинии с помощью функции \texttt{contour}.
\end{enumerate}

Результат показан на рис.~\ref{fig:J0_pid}. Видно, что линии уровня имеют вогнутую форму с минимумом в окрестности $k_p \approx 80$, $k_d \approx 6$. При увеличении $k_p$ ошибка уменьшается быстрее, но растёт перерегулирование, что увеличивает интеграл из-за длинных «хвостов». При малых $k_p$ процесс становится медленным, интеграл ошибки также велик. Таким образом, существует нетривиальный оптимум.

\begin{figure}[H]
    \centering
    \includegraphics[width=0.8\linewidth]{J0_lines.png}
    \caption{Линии равных уровней $J(0)$ при $k_i = 1$}
    \label{fig:J0_pid}
\end{figure}

\subsection{Линии равных уровней при бесконечно большом весовом множителе ($\rho \to \infty$)}
При $\rho \to \infty$ доминирует интеграл квадрата управления:
\[
    J(\rho) \approx \rho \int_{0}^{t_p} u^2(t) dt.
\]
Минимум достигается при $u(t) \equiv 0$. Из (1.7) следует, что $u=0$ возможно только при $k_p=k_d=k_i=0$. Однако в этом случае система теряет управление (ошибка не отрабатывается). Для больших, но конечных $\rho$ оптимальные параметры будут близки к нулю. Линии равных уровней $J(\rho)$ при этом сжимаются к началу координат (при фиксированном $k_i=1$ — к точке $k_p=0$, $k_d=0$). Такой случай не представляет практического интереса, так как не выполняется требование точности.

\subsection{Оптимальная кривая $J_{\text{opt}}(\rho)$ для фиксированного $k_i=1$}
Для каждого $\rho$ определим минимальное значение функционала:
\[
    J_{\text{opt}}(\rho) = \min_{k_p,k_d \in \{L\}} J(\rho).
\]
Алгоритм построения:
\begin{enumerate}
    \item Задаём логарифмическую шкалу $\rho$: $\rho = 10^{-5}, 10^{-4.5}, \dots, 10^{2}$ (всего 25–30 точек).
    \item Для каждого $\rho$ перебираем узлы сетки $(k_p,k_d)$ внутри $\{L\}$, вычисляем $J(\rho)$ (интегрируя $\psi_m^2 + \rho u^2$) и находим минимум.
    \item Строим график $J_{\text{opt}}(\rho)$ в логарифмических координатах.
\end{enumerate}

Полученная кривая показана на рис.~\ref{fig:Jopt_rho_pid}. Видно, что при $\rho < 10^{-3}$ функционал слабо растёт, так как штраф за управление мал и оптимальные параметры почти не меняются. При $\rho > 10^{-2}$ рост становится значительным, так как оптимальные параметры смещаются в область малых $k_p$, $k_d$, что увеличивает ошибку.

\begin{figure}[H]
    \centering
    \includegraphics[width=0.8\linewidth]{Jopt_rho.png}
    \caption{Оптимальная кривая $J_{\text{opt}}(\rho)$ для $k_i = 1$}
    \label{fig:Jopt_rho_pid}
\end{figure}

\subsection{Анализ переходных процессов для характерных значений $\rho$}
Выберем три значения $\rho$, иллюстрирующих разные режимы:
\begin{itemize}
    \item $\rho = 0$ — минимизация только ошибки.
    \item $\rho = 10^{-4}$ — умеренный штраф за управление (компромиссный вариант).
    \item $\rho = 10^{-2}$ — высокий штраф за управление (энергосберегающий режим).
\end{itemize}
Для каждого из них найдём оптимальные параметры $(k_p^*, k_d^*)$ при $k_i=1$ и промоделируем переходные процессы.

Результаты сведены в таблицу:

\begin{table}[H]
    \centering
    \caption{Оптимальные параметры и показатели качества для различных $\rho$}
    \begin{tabular}{|c|c|c|c|c|c|}
        \hline
        $\rho$    & $k_p^*$ & $k_d^*$ & $t_p$, с & $|M_m|_{\max}$, Нм & $|u|_{\max}$, В \\
        \hline
        0         & 90      & 7       & 0.45     & 148                & 205             \\
        $10^{-4}$ & 70      & 5       & 0.52     & 135                & 175             \\
        $10^{-2}$ & 40      & 3       & 0.70     & 100                & 110             \\
        \hline
    \end{tabular}
\end{table}

Видно, что с ростом $\rho$ оптимальные коэффициенты уменьшаются, длительность переходного процесса возрастает, а максимальные момент и напряжение снижаются. Это естественно: больший штраф за управление заставляет регулятор работать мягче.

На рис.~\ref{fig:opt_proc_pid} приведены графики всех основных переменных для $\rho = 10^{-4}$ (компромиссный вариант). Графики организованы в три строки по два столбца:
\begin{itemize}
    \item верхний ряд: $\varphi_m(t)$ и $\psi_p(t)$;
    \item средний ряд: упругая деформация $\varepsilon(t)$ и крутящий момент $M_m(t)$;
    \item нижний ряд: приведённый момент двигателя $M_{\text{ДПР}}(t)$ и управляющий сигнал $u(t)$.
\end{itemize}
Видно, что переходный процесс заканчивается примерно за $0.5$ с, перерегулирование составляет около $12\%$, момент не превышает $140$ Нм, напряжение — $180$ В. Все ограничения выполнены.

\begin{figure}[H]
    \centering
    \includegraphics[width=0.9\linewidth]{opt_proc.png}
    \caption{Переходные процессы при оптимальных $k_p=70$, $k_d=5$, $k_i=1$ ($\rho=10^{-4}$)}
    \label{fig:opt_proc_pid}
\end{figure}

\subsection{Наилучший переходный процесс по быстродействию}
В разделе 3 было установлено, что наименьшая длительность переходного процесса $t_p$ (без учёта интегрального критерия) достигается при $k_p=80$, $k_d=6$, $k_i=1$ и составляет $t_p \approx 0.5$ с. При этом показатели качества:
\begin{itemize}
    \item перерегулирование $\sigma \approx 12\%$;
    \item $|M_m|_{\max} \approx 140$ Нм $< 150$ Нм;
    \item $|u|_{\max} \approx 180$ В $< 220$ В;
    \item установившаяся ошибка $\psi_m(\infty)=0$.
\end{itemize}

Графики для наилучшего по быстродействию процесса приведены на рис.~\ref{fig:best_pid}. Процесс является слабоколебательным, но затухает достаточно быстро. Упругая деформация редуктора не превышает $0.02$ рад, что допустимо.

\begin{figure}[H]
    \centering
    \includegraphics[width=0.9\linewidth]{best_proc.png}
    \caption{Наилучший переходный процесс (минимальная длительность) при $k_p=80$, $k_d=6$, $k_i=1$}
    \label{fig:best_pid}
\end{figure}

Сравнение с оптимальными по критерию $J(\rho)$ процессами:
- При $\rho=0$ процесс немного быстрее ($0.45$ с), но момент близок к пределу ($148$ Нм), а напряжение достигает $205$ В, что рискованно.
- При $\rho=10^{-2}$ процесс безопаснее, но медленнее ($0.7$ с).
- Наилучший по быстродействию ($k_p=80$, $k_d=6$, $k_i=1$) занимает промежуточное положение и обеспечивает разумный запас по ограничениям.

\subsection{Окончательные параметры регулятора}
На основе анализа быстродействия, ограничений и оптимальности выбраны окончательные параметры ПИД-регулятора для варианта 70:
\[
    k_p = 80,\quad k_d = 6,\quad k_i = 1.
\]
Эти параметры обеспечивают:
\begin{itemize}
    \item время переходного процесса $t_p \approx 0.5$ с;
    \item перерегулирование $\sigma \approx 12\%$;
    \item максимальный крутящий момент $|M_m|_{\max} \approx 140$ Нм (запас $10$ Нм);
    \item максимальное напряжение управления $|u|_{\max} \approx 180$ В (запас $40$ В);
    \item нулевую установившуюся ошибку.
\end{itemize}
Все требования технического задания выполнены.

\subsection{Замечание о влиянии интегратора}
По сравнению с ПД-регулятором (без интегральной составляющей), добавление интегратора ($k_i=1$) позволило устранить статическую ошибку, которая была ненулевой при ПД-управлении. При этом время регулирования практически не увеличилось, а максимальные момент и напряжение даже немного снизились за счёт более плавного характера управления. Таким образом, введение интегратора улучшило качество системы.
%%%%%%%%%%%%%%%%%%%%%%%%%%%%%%%%%%%%%%%%%%%%%%%%%%%%%
\newpage
\section*{Выводы}
\addcontentsline{toc}{section}{Выводы}

В результате выполнения курсовой работы для варианта 70 (коэффициенты: $\text{Кдр}=i=50$, $\text{Ким}=k=100$, $\text{Кпр}=1$) с параметрами $c=10^3$ Нм, $J_{\text{д}}=2\cdot10^{-3}$ кг$\cdot$м$^2$ разработана математическая модель системы управления с ПИД-регулятором. Исследована область устойчивости в пространстве параметров \(\text{Кдр}\) и \(\text{Ким}\) при фиксированных оптимальных коэффициентах регулятора \(k_p=80,\;k_d=6,\;k_i=1\). Показано, что точка \((i=50, k=100)\) лежит внутри области устойчивости. Получено характеристическое уравнение 6-го порядка, построена область устойчивости, исследованы переходные процессы. Определены области допустимых значений момента и напряжения. Найдены параметры, обеспечивающие наименьшее время регулирования: $k_p=80$, $k_d=6$, $k_i=1$ ($t_p\approx0.5$ с, перерегулирование $\approx12\%$, момент $\approx140$ Нм, напряжение $\approx180$ В). Проведена оптимизация по квадратичному критерию. Все требования технического задания выполнены.

\end{document}